\begin{abstract}

Modern web browsers enforce isolation between websites so that a malicious webpage should not be able to directly read the contents of another origin. However, side-channel attacks can sometimes infer information indirectly through shared browser and hardware resources. In this project, we propose to study and reproduce a pixel-stealing attack in a controlled environment. We will create a malicious webpage and a separate victim webpage containing known visual information, then investigate whether rendering-related side channels can be used to recover information from the victim page while the user remains on the malicious site. Our primary goal is to determine whether such an attack can be reproduced, how accurately information can be recovered, and how long recovery takes. As a stretch goal, we will investigate whether partial recovery progress can be saved and resumed if the victim later revisits the malicious webpage.

\end{abstract}
